\documentclass[a4paper]{article}
\usepackage[pdftex]{graphicx}
\usepackage[utf8]{inputenc}
\usepackage{enumerate}
\usepackage{siunitx}
\sisetup{locale=DE} 
\usepackage{href-ul}
\hypersetup{
	colorlinks=true,
	linkcolor=blue,
	urlcolor=blue}
\usepackage{icomma}
\usepackage{amssymb}
\usepackage{geometry}
\geometry{a4paper, top=15mm, left=15mm, right=15mm, bottom=15mm,
	headsep=10mm, footskip=12mm}

\begin{document}
{\bf Elektrizitätslehre}
\begin{enumerate}[1.]
%Aufgabe 2
{\bf \item Der zerstreute  ,,Physiker" Dr. Keine-Ahnung behauptet:}

\begin{enumerate}[a)]
\item ,,Den Südpol eines eingeschalteten Elektromagneten erkennt man daran, dass der Eisenkern an diesem Ende rot eingefärbt ist."
\item ,,So erkennt man den Minuspol einer Gleichstromquelle: In einem unverzweigten Stromkreis mit zwei baugleichen Glühlampen leuchtet diejenige heller, die näher am Pluspol der Stromquelle ist."
\end{enumerate}
Was ist an diesen Behauptungen falsch? Beschreibe das richtige Vorgehen.

{\bf \item In das Innere einer langen Spule legt man hintereinander längs zur Spulenachse zwei dünne, zylindrische Eisenstäbe.} 
\begin{enumerate}[a)]
\item Erläutere und begründe, was mit den Eisenstäben geschieht, wenn durch die Spule Strom fließt.
\item Nun legt man man die Stäbe nebeneinander in die Spule und schaltet den Strom ein. Was ist jetzt zu beobachten?
\end{enumerate}

{\bf \item Elektrische Schaltungen: In einem Stromkreis mit einer Stromquelle und einer Glühlampe soll 
$\dots$}
\begin{enumerate}[a)]
\item $\dots$ die Stromstärke  $\dots$ 
\item $\dots$ die Betriebsspannung  $\dots$ 
\end{enumerate}
gemessen werden. Zeichne zu beiden Fällen je eine Schaltskizze


%Aufgabe 3
{\bf \item Aufgabe}
\begin{enumerate}[a)]
\item Zeichne einen Stromkreis mit Batterie und zwei in Reihe geschalteten Glühbirnen. 
\item Die Stromstärke an beiden Glühbirnen soll gemessen werden. Gib an, wieso dazu ein Messgerät genügt und zeichne eine neue Schaltung mit dem Messgerät.
\item Nun sollen die Spannungen an beiden Glühbirnen gemessen werden. Zeichne eine Schaltung, mit der dies möglich ist. 
\end{enumerate}

%Aufgabe 4
{\bf \item Aufgabe}
\begin{enumerate}[a)]
\item An einem Akku mit der Aufschrift $\SI{12}{\volt} \quad \SI{3000}{\milli\ampere\hour}$ wird ein Notebook mit der Leistung $\SI{24}{\watt}$ angeschlossen. Wie lange kann das Notebook mit diesem Akku betrieben werden?
\item Wieviel Höhenmeter könnte ein mit diesem Akku betriebenes E-Bike einen Berg hinauffahren, wenn es mit Fahrer $ \SI{105}{\kilogram}$ wiegt und man annimmt, dass diese Energie vollständig in Lageenergie umgewandelt wird?
\end{enumerate}

\end{enumerate} 

\fbox{
	\begin{minipage}{0.5\textwidth}
		Zur Lösung bitte \href{https://www.okuyakl.de/phys/p9eleL403/ll403.pdf}{hier klicken} oder den QR-Code scannen.\\
	Weitere Arbeitsblätter gibt es unter 
	
	\href{https://www.okuyakl.de}{www.okuyakl.de}
	\end{minipage}
	\hfill
	\begin{minipage}{0.4\textwidth}
		\includegraphics[width=1.5 cm]{../../viecher/zwe03}
		\includegraphics[width=3 cm]{qr403}
		\includegraphics[width=2 cm]{../../viecher/afanticon1}
		
	\end{minipage}}

\end{document}